\documentclass[11pt,a4paper]{article}
\usepackage[margin=1in]{geometry}
\usepackage{booktabs}
\usepackage{float}
\usepackage{makecell}
\usepackage{placeins}
\usepackage{caption}
\usepackage{amsmath,amssymb}
\usepackage{hyperref}
\title{\textbf{Section 3.1 IM-OLS Cointegration Estimation Matrix: U.S. NFC Transformation Elasticity (1929--2024)}}
\author{\textbf{Department of Economics, UMass Amherst}}
\date{August 2026}
\begin{document}
\maketitle

\section*{Data Provenance and Estimation Summary}
This report presents the empirical cointegration results for Section 3.1 of Chapter 2. All models are estimated via Cointegrating Multivariate Polynomial Regression (CMPR) IM-OLS (Vogelsang et al. 2024) over the full NIPA sample window (1931--2024, $N=94$ observations after listwise deletion for the lag-1 conditioning variables $\omega_{t-1}$ and $\iota_{t-1}$). Point estimates are tuning-parameter free; coefficient and cointegration-test significance stars use Vogelsang \& Wagner (2014) fixed-$b$ critical values (Bartlett kernel, Newey-West bandwidth), interpolated by each model's bandwidth ratio and number of stochastic regressors -- see the significance note under each table. Distributional conditioning variables ($\omega_{t-1}$ and $\iota_{t-1}$) are centered at their 1929 initial values, preserving $\hat{\theta}_0$ as the capital elasticity at the initial state. VIF statistics are computed on the level regressors space. To ensure robustness against structural instability, each specification dynamically incorporates step level-shift dummies $D^{\text{step}}_{T_b}$ for structural breaks identified by a first-pass Bai-Perron residual audit.

\FloatBarrier
\section{Specification Set 1: Double Interactions ($k^{\text{cap}} \cdot \omega_{t-1}$)}
\FloatBarrier
\begin{table}[H]
\centering
\caption{Specification Set 1: Double Interaction IM-OLS Cointegration Estimates (1929--2024)}
\label{tab:spec_set_1}
\scriptsize
\begin{tabular}{p{3.8cm}p{1.35cm}p{1.35cm}p{1.35cm}p{1.35cm}p{1.35cm}p{1.35cm}}
\toprule
\multicolumn{7}{l}{\textbf{Dependent Variable: $y_t^{\text{NFC}}$ (Log NFC Real GVA)}} \\
\midrule
Model & M1.1 & M1.2 & M1.3 & M1.4 & M1.5 & M1.6  \\
Capital Stock & $k^{\text{cap}}_1$ & $k^{\text{cap}}_2$ & $k^{\text{cap}}_1$ & $k^{\text{cap}}_2$ & $k^{\text{cap}}_1$ & $k^{\text{cap}}_2$  \\
Demand Control & None & None & None & None & None & None  \\
\midrule
$\hat{\alpha}$ (Intercept)  & 8.97 & 6.81 & 9.19 & 8.75 & 5.44 & 4.84  \\
 & {\scriptsize (0.89)} & {\scriptsize (0.83)} & {\scriptsize (0.98)} & {\scriptsize (0.90)} & {\scriptsize (1.23)} & {\scriptsize (1.24)}  \\
Step Shifts ($D^{	ext{step}}_{T_b}$)  &   \
$\hat{\theta}_0$ (Capital Elasticity)  & 0.34 & 0.52 & 0.33 & 0.35 & 0.67 & 0.70  \\
 & {\scriptsize (0.08)} & {\scriptsize (0.07)} & {\scriptsize (0.09)} & {\scriptsize (0.08)} & {\scriptsize (0.11)} & {\scriptsize (0.10)}  \\
$\hat{\theta}_1$ (Interaction)  & -0.05 & -0.07 & -0.07 & -0.04 & -0.19 & -0.18  \\
 & {\scriptsize (0.15)} & {\scriptsize (0.10)} & {\scriptsize (0.16)} & {\scriptsize (0.14)} & {\scriptsize (0.13)} & {\scriptsize (0.13)}  \\
\midrule
\multicolumn{7}{l}{\underline{Collinearity Diagnostics (VIF)}} \\
VIF ($k^{\text{cap}}$) & 1.38 & 1.40 & 1.54 & 1.56 & 8.04 & 8.14  \\
VIF (Interaction) & 1.38 & 1.40 & 1.54 & 1.56 & 8.04 & 8.14  \\
\midrule
\multicolumn{7}{l}{\underline{Cointegration Test Triad}} \\
\makecell[l]{Resid. ADF ($t_{\text{ADF}}$) \\ {\scriptsize [$H_0$: No Coint.]}} & -3.00 & -3.52$^{*}$ & -3.15 & -2.95 & -2.78 & -2.86  \\
\makecell[l]{Shin LM ($C_{\text{Shin}}$) \\ {\scriptsize [$H_0$: Coint.]}} & 0.56$^{***}$ & 0.92$^{***}$ & 0.45$^{***}$ & 0.60$^{***}$ & 0.20$^{***}$ & 0.20$^{***}$  \\
\makecell[l]{VW Type 2 ($S_{\text{Type 2}}$) \\ {\scriptsize [$H_0$: Coint.]}} & 1835.00$^{***}$ & 719.49$^{***}$ & 1794.32$^{***}$ & 2306.76$^{***}$ & 2169.68$^{***}$ & 1404.04$^{***}$  \\
\midrule
$R^2_{\text{sum}}$ & 1.00 & 1.00 & 1.00 & 1.00 & 1.00 & 1.00  \\
Sample Size ($T$) & 94 & 94 & 94 & 94 & 94 & 94  \\
\bottomrule
\multicolumn{7}{p{13.2cm}}{\footnotesize Note: Estimated via Cointegrating Multivariate Polynomial Regression (CMPR) IM-OLS (Vogelsang et al. 2024) over 1929--2024. Standard errors in parentheses. Significance: $^*p<0.10, ^{**}p<0.05, ^{***}p<0.01$. Coefficient rows ($\hat{\alpha}$, $\hat{\gamma}^{\text{step}}$, $\hat{\theta}_0$, $\hat{\theta}_1$, $\hat{\beta}_{\text{cm}}$) and the Residual ADF and VW Type 2 rows use Vogelsang \& Wagner (2014) fixed-$b$ critical values, interpolated by each model's bandwidth ratio $b=\text{bandwidth}/T$ (Bartlett kernel, Newey-West bandwidth) and its number of stochastic regressors $k$ -- NOT a single generic $t$/normal cutoff. Resid. ADF stars instead follow MacKinnon (2010) Engle-Granger critical values, indexed by the total number of $I(1)$ variables; the test is lag-augmented with the order selected by AIC (max 8), which is required because the IM-OLS residuals are strongly autocorrelated -- an unaugmented Dickey-Fuller is biased toward non-rejection and was the source of the superseded v2 results. Shin LM stars use Shin's (1994) critical values (10\%/5\%/1\% = 0.095/0.121/0.184); these have not been independently re-verified as regressor-count-specific for this report's $k=2$/$k=3$ mix and should be treated as provisional (see the audit memo referenced in the repository). \textbf{Star reversal warning:} For Residual ADF, stars indicate rejection of No Cointegration (supporting cointegration). For Shin LM and VW Type 2, stars indicate rejection of Cointegration. Distributional variables centered at 1929 initial state ($t_0$). Step dummies $D^{\text{step}}_{T_b} = \mathbb{I}(t \ge T_b)$ represent persistent level shifts from structural breaks identified via a first-pass Bai-Perron residual audit (15\% trimming), distinguished from impulse outlier controls $D^{\text{imp}}_{T_b} = \mathbb{I}(t = T_b)$.}
\end{tabular}
\end{table}


\FloatBarrier
\section{Specification Set 2: Triple Interactions ($k^{\text{cap}} \cdot \omega_{t-1} \cdot \iota_{t-1}$)}
\FloatBarrier
\begin{table}[H]
\centering
\caption{Specification Set 2 (Part A: Baseline NFC \& Corporate GVA Wage Shares): Triple Interaction IM-OLS Cointegration Estimates}
\label{tab:spec_set_2_a}
\scriptsize
\begin{tabular}{p{3.8cm}p{1.35cm}p{1.35cm}p{1.35cm}p{1.35cm}p{1.35cm}p{1.35cm}}
\toprule
\multicolumn{7}{l}{\textbf{Dependent Variable: $y_t^{\text{NFC}}$ (Log NFC Real GVA)}} \\
\midrule
Model & M2.1 & M2.2 & M2.3 & M2.4 & M2.5 & M2.6  \\
Capital Stock & $k^{\text{cap}}_1$ & $k^{\text{cap}}_1$ & $k^{\text{cap}}_2$ & $k^{\text{cap}}_2$ & $k^{\text{cap}}_1$ & $k^{\text{cap}}_1$  \\
Demand Control & None & NetC Share & None & NetC Share & None & NetC Share  \\
\midrule
$\hat{\alpha}$ (Intercept)  & 8.18 & 9.41 & 8.08 & 7.22 & 7.76 & 9.36  \\
 & {\scriptsize (1.06)} & {\scriptsize (0.88)} & {\scriptsize (1.10)} & {\scriptsize (0.88)} & {\scriptsize (1.16)} & {\scriptsize (0.89)}  \\
Step Shifts ($D^{	ext{step}}_{T_b}$)  &   \
$\hat{\theta}_0$ (Capital Elasticity)  & 0.42 & 0.36 & 0.41 & 0.50 & 0.45 & 0.36  \\
 & {\scriptsize (0.09)} & {\scriptsize (0.07)} & {\scriptsize (0.09)} & {\scriptsize (0.06)} & {\scriptsize (0.10)} & {\scriptsize (0.07)}  \\
$\hat{\theta}_1$ (Interaction)  & -0.19 & -0.09 & -0.17 & -0.46 & -0.16 & -0.07  \\
 & {\scriptsize (0.09)} & {\scriptsize (0.08)} & {\scriptsize (0.10)} & {\scriptsize (0.13)} & {\scriptsize (0.10)} & {\scriptsize (0.08)}  \\
$\hat{\beta}_{\text{cm}}$ ($\text{NetC}^{\text{share}}$)  & --- & -0.93 & --- & -0.25 & --- & -0.92  \\
 &  & {\scriptsize (0.62)} &  & {\scriptsize (0.60)} &  & {\scriptsize (0.63)}  \\
\midrule
\multicolumn{7}{l}{\underline{Collinearity Diagnostics (VIF)}} \\
VIF ($k^{\text{cap}}$) & 3.57 & 4.41 & 3.69 & 4.61 & 3.86 & 4.64  \\
VIF (Interaction) & 3.57 & 3.58 & 3.69 & 3.71 & 3.86 & 3.87  \\
\midrule
\multicolumn{7}{l}{\underline{Cointegration Test Triad}} \\
\makecell[l]{Resid. ADF ($t_{\text{ADF}}$) \\ {\scriptsize [$H_0$: No Coint.]}} & -4.00$^{**}$ & -2.64 & -3.00 & -3.21 & -3.93$^{**}$ & -2.62  \\
\makecell[l]{Shin LM ($C_{\text{Shin}}$) \\ {\scriptsize [$H_0$: Coint.]}} & 0.66$^{***}$ & 0.83$^{***}$ & 0.71$^{***}$ & 0.41$^{***}$ & 0.72$^{***}$ & 0.83$^{***}$  \\
\makecell[l]{VW Type 2 ($S_{\text{Type 2}}$) \\ {\scriptsize [$H_0$: Coint.]}} & 276.16$^{***}$ & --- & 359.15$^{***}$ & --- & 140.50$^{***}$ & ---  \\
\midrule
$R^2_{\text{sum}}$ & 1.00 & 1.00 & 1.00 & 1.00 & 1.00 & 1.00  \\
Sample Size ($T$) & 94 & 94 & 94 & 94 & 94 & 94  \\
\bottomrule
\multicolumn{7}{p{13.2cm}}{\footnotesize Note: Estimated via Cointegrating Multivariate Polynomial Regression (CMPR) IM-OLS (Vogelsang et al. 2024) over 1929--2024. Standard errors in parentheses. Significance: $^*p<0.10, ^{**}p<0.05, ^{***}p<0.01$. Coefficient rows ($\hat{\alpha}$, $\hat{\gamma}^{\text{step}}$, $\hat{\theta}_0$, $\hat{\theta}_1$, $\hat{\beta}_{\text{cm}}$) and the Residual ADF and VW Type 2 rows use Vogelsang \& Wagner (2014) fixed-$b$ critical values, interpolated by each model's bandwidth ratio $b=\text{bandwidth}/T$ (Bartlett kernel, Newey-West bandwidth) and its number of stochastic regressors $k$ -- NOT a single generic $t$/normal cutoff. Resid. ADF stars instead follow MacKinnon (2010) Engle-Granger critical values, indexed by the total number of $I(1)$ variables; the test is lag-augmented with the order selected by AIC (max 8), which is required because the IM-OLS residuals are strongly autocorrelated -- an unaugmented Dickey-Fuller is biased toward non-rejection and was the source of the superseded v2 results. Shin LM stars use Shin's (1994) critical values (10\%/5\%/1\% = 0.095/0.121/0.184); these have not been independently re-verified as regressor-count-specific for this report's $k=2$/$k=3$ mix and should be treated as provisional (see the audit memo referenced in the repository). \textbf{Star reversal warning:} For Residual ADF, stars indicate rejection of No Cointegration (supporting cointegration). For Shin LM and VW Type 2, stars indicate rejection of Cointegration. Distributional variables centered at 1929 initial state ($t_0$). Step dummies $D^{\text{step}}_{T_b} = \mathbb{I}(t \ge T_b)$ represent persistent level shifts from structural breaks identified via a first-pass Bai-Perron residual audit (15\% trimming), distinguished from impulse outlier controls $D^{\text{imp}}_{T_b} = \mathbb{I}(t = T_b)$.}
\end{tabular}
\end{table}


\FloatBarrier
\begin{table}[H]
\centering
\caption{Specification Set 2 (Part B: Shaikh-Tonak Tier 2 Wage Share): Triple Interaction IM-OLS Cointegration Estimates}
\label{tab:spec_set_2_b}
\scriptsize
\begin{tabular}{p{3.8cm}p{1.35cm}p{1.35cm}p{1.35cm}p{1.35cm}p{1.35cm}p{1.35cm}}
\toprule
\multicolumn{7}{l}{\textbf{Dependent Variable: $y_t^{\text{NFC}}$ (Log NFC Real GVA)}} \\
\midrule
Model & M2.7 & M2.8 & M2.9 & M2.10 & M2.11 & M2.12  \\
Capital Stock & $k^{\text{cap}}_2$ & $k^{\text{cap}}_2$ & $k^{\text{cap}}_1$ & $k^{\text{cap}}_1$ & $k^{\text{cap}}_2$ & $k^{\text{cap}}_2$  \\
Demand Control & None & NetC Share & None & NetC Share & None & NetC Share  \\
\midrule
$\hat{\alpha}$ (Intercept)  & 7.93 & 8.60 & 8.68 & 9.49 & 8.44 & 8.98  \\
 & {\scriptsize (1.13)} & {\scriptsize (1.00)} & {\scriptsize (0.70)} & {\scriptsize (0.81)} & {\scriptsize (0.73)} & {\scriptsize (0.88)}  \\
Step Shifts ($D^{	ext{step}}_{T_b}$)  &   \
$\hat{\theta}_0$ (Capital Elasticity)  & 0.43 & 0.41 & 0.39 & 0.43 & 0.40 & 0.45  \\
 & {\scriptsize (0.10)} & {\scriptsize (0.08)} & {\scriptsize (0.06)} & {\scriptsize (0.06)} & {\scriptsize (0.06)} & {\scriptsize (0.07)}  \\
$\hat{\theta}_1$ (Interaction)  & -0.15 & -0.01 & -0.38 & -0.30 & -0.36 & -0.29  \\
 & {\scriptsize (0.10)} & {\scriptsize (0.09)} & {\scriptsize (0.11)} & {\scriptsize (0.07)} & {\scriptsize (0.11)} & {\scriptsize (0.07)}  \\
$\hat{\beta}_{\text{cm}}$ ($\text{NetC}^{\text{share}}$)  & --- & -0.89 & --- & -2.12 & --- & -1.99  \\
 &  & {\scriptsize (0.63)} &  & {\scriptsize (0.39)} &  & {\scriptsize (0.40)}  \\
\midrule
\multicolumn{7}{l}{\underline{Collinearity Diagnostics (VIF)}} \\
VIF ($k^{\text{cap}}$) & 4.00 & 4.84 & 14.01 & 14.05 & 14.46 & 14.48  \\
VIF (Interaction) & 4.00 & 4.01 & 14.01 & 15.20 & 14.46 & 15.56  \\
\midrule
\multicolumn{7}{l}{\underline{Cointegration Test Triad}} \\
\makecell[l]{Resid. ADF ($t_{\text{ADF}}$) \\ {\scriptsize [$H_0$: No Coint.]}} & -2.98 & -2.50 & -2.90 & -3.06 & -2.85 & -3.09  \\
\makecell[l]{Shin LM ($C_{\text{Shin}}$) \\ {\scriptsize [$H_0$: Coint.]}} & 0.76$^{***}$ & 0.90$^{***}$ & 0.51$^{***}$ & 0.38$^{***}$ & 0.52$^{***}$ & 0.39$^{***}$  \\
\makecell[l]{VW Type 2 ($S_{\text{Type 2}}$) \\ {\scriptsize [$H_0$: Coint.]}} & 223.87$^{***}$ & --- & 3570.21$^{***}$ & --- & 3285.32$^{***}$ & ---  \\
\midrule
$R^2_{\text{sum}}$ & 1.00 & 1.00 & 1.00 & 1.00 & 1.00 & 1.00  \\
Sample Size ($T$) & 94 & 94 & 94 & 94 & 94 & 94  \\
\bottomrule
\multicolumn{7}{p{13.2cm}}{\footnotesize Note: Estimated via Cointegrating Multivariate Polynomial Regression (CMPR) IM-OLS (Vogelsang et al. 2024) over 1929--2024. Standard errors in parentheses. Significance: $^*p<0.10, ^{**}p<0.05, ^{***}p<0.01$. Coefficient rows ($\hat{\alpha}$, $\hat{\gamma}^{\text{step}}$, $\hat{\theta}_0$, $\hat{\theta}_1$, $\hat{\beta}_{\text{cm}}$) and the Residual ADF and VW Type 2 rows use Vogelsang \& Wagner (2014) fixed-$b$ critical values, interpolated by each model's bandwidth ratio $b=\text{bandwidth}/T$ (Bartlett kernel, Newey-West bandwidth) and its number of stochastic regressors $k$ -- NOT a single generic $t$/normal cutoff. Resid. ADF stars instead follow MacKinnon (2010) Engle-Granger critical values, indexed by the total number of $I(1)$ variables; the test is lag-augmented with the order selected by AIC (max 8), which is required because the IM-OLS residuals are strongly autocorrelated -- an unaugmented Dickey-Fuller is biased toward non-rejection and was the source of the superseded v2 results. Shin LM stars use Shin's (1994) critical values (10\%/5\%/1\% = 0.095/0.121/0.184); these have not been independently re-verified as regressor-count-specific for this report's $k=2$/$k=3$ mix and should be treated as provisional (see the audit memo referenced in the repository). \textbf{Star reversal warning:} For Residual ADF, stars indicate rejection of No Cointegration (supporting cointegration). For Shin LM and VW Type 2, stars indicate rejection of Cointegration. Distributional variables centered at 1929 initial state ($t_0$). Step dummies $D^{\text{step}}_{T_b} = \mathbb{I}(t \ge T_b)$ represent persistent level shifts from structural breaks identified via a first-pass Bai-Perron residual audit (15\% trimming), distinguished from impulse outlier controls $D^{\text{imp}}_{T_b} = \mathbb{I}(t = T_b)$.}
\end{tabular}
\end{table}


\end{document}
